%------------------------
% Resume Template
% Author : Anubhav Singh
% Github : https://github.com/xprilion
% License : MIT
%------------------------

\documentclass[a4paper]{article}

\RequirePackage{fix-cm}
\usepackage{latexsym}
\usepackage[empty]{fullpage}
\usepackage{titlesec}
\usepackage{marvosym}
\usepackage[usenames,dvipsnames]{color}
\usepackage{verbatim}
\usepackage{enumitem}
\usepackage[pdftex,hidelinks]{hyperref}
\usepackage{fancyhdr}

\pagestyle{fancy}
\fancyhf{}
\fancyfoot{}
\renewcommand{\headrulewidth}{0pt}
\renewcommand{\footrulewidth}{0pt}

\usepackage[T1]{fontenc}
\usepackage{lmodern}
\input{glyphtounicode}
\pdfgentounicode=1
\usepackage{microtype}
\DisableLigatures[f]{encoding = *, family = *}

% Margins
\addtolength{\oddsidemargin}{-0.530in}
\addtolength{\evensidemargin}{-0.375in}
\addtolength{\textwidth}{1in}
\addtolength{\topmargin}{-0.7in}
\addtolength{\textheight}{1.35in}

\urlstyle{rm}
\raggedbottom
\raggedright
\setlength{\tabcolsep}{0in}

% Section headings
\titleformat{\section}{\vspace{-4pt}\scshape\raggedright\large}{}{0em}{}[\color{black}\titlerule\vspace{-5pt}]

%-------------------------
% Custom commands
\newcommand{\resumeItem}[2]{\item\fontsize{9.8}{11.9}\selectfont{\textbf{#1}{: #2}}}
\newcommand{\bul}[1]{\item\fontsize{9.8}{11.9}\selectfont{#1}}
\newcommand{\resumeSubheading}[4]{
  \vspace{1pt}\item
    \begin{tabular*}{0.985\textwidth}{l@{\extracolsep{\fill}}r}
      \textbf{#1} & #2 \\
      \textit{#3} & \textit{#4} \\
    \end{tabular*}\vspace{-1pt}
}
\newcommand{\resumeSubItem}[2]{\resumeItem{#1}{#2}\vspace{-1pt}}
\newcommand{\resumeSubHeadingListStart}{\begin{itemize}[leftmargin=0pt,label={},itemsep=2pt,topsep=3pt]}
\newcommand{\resumeSubHeadingListEnd}{\end{itemize}}
\newcommand{\resumeItemListStart}{\begin{itemize}[leftmargin=12pt,itemsep=1.6pt,topsep=2pt,parsep=0pt]}
\newcommand{\resumeItemListEnd}{\end{itemize}\vspace{-3pt}}

%-----------------------------
\begin{document}
\fontfamily{cmss}\selectfont

%---------- HEADING ----------
\begin{tabular*}{\textwidth}{l@{\extracolsep{\fill}}r}
  \textbf{{\LARGE Rahul Muralidhara Koushik}} & San Francisco Bay Area \\
  Email: \href{mailto:rkoushik@asu.edu}{rkoushik@asu.edu} & Mobile:~+1-480-791-6597 \\
  \href{https://github.com/Fluff18}{GitHub: github.com/Fluff18} \;|\; \href{https://fluff18.github.io/v3/}{Portfolio: fluff18.github.io/v3} & LinkedIn: \href{https://www.linkedin.com/in/rahulkoushik/}{linkedin.com/in/rahulkoushik} \\
\end{tabular*}

\vspace{2pt}
{\fontsize{9.8}{11.9}\selectfont \textbf{AI Engineer \textbar{} LLM Agents, Evals, Inference.} Builds generative AI (GenAI) and large language model (LLM) agents end to end: the user problem, the agent loop, the memory layer and the evals that decide whether an agent is ready to ship. Open to AI Engineer and AI Product roles in the San Francisco Bay Area.}

%---------- EXPERIENCE ----------
\section{Experience}
\resumeSubHeadingListStart
  \resumeSubheading{Hewlett Packard Enterprise, Storage AI Team}{San Jose, US}{AI Engineer}{Feb 2026 -- Present}
  \resumeItemListStart
    \bul{Architected the production \textbf{multi-agent LLM harness} with \textbf{guardrails} that every team agent runs on; its \textbf{Model Context Protocol (MCP)} tool server handles \textbf{6,300+ tool calls a week} across HPE Storage and backs the storage experts' Copilot.}
    \bul{Built the \textbf{memory layer} for persistent agents: episodic memory across runs, \textbf{embeddings and vector search} on pgvector, and a \textbf{context graph} linking \textbf{13,005} mined bugs and \textbf{5,600} grounded analyses by component, symptom and fix, built in nine days.}
    \bul{Cut agent latency and cost per run with \textbf{inference optimization}: \textbf{prompt caching} and \textbf{KV-cache reuse} across agent steps, \textbf{cache-coherent} context updates, and \textbf{prefill/decode-aware} serving of \textbf{Qwen} models (27B, 35B) on \textbf{vLLM}.}
    \bul{Took the \textbf{agent evaluation platform} from idea to \textbf{three production deployments in four weeks}: rubric-based evals scored against expert feedback, \textbf{LLM-as-a-judge}, and a \textbf{tracing dashboard} showing every step, tool call and failure of each agent run. It is the \textbf{release gate} every agent must pass.}
    \bul{Shipped \textbf{agentic alert-triage and root-cause-analysis (RCA)} systems to production, with \textbf{self-improving loops} where agents mine their own failures, propose patches and re-run against the pass bar.}
    \bul{Led the self-learning agent experiments, including reporting an honest \textbf{zero-lift} result after catching data leakage in the training signal.}
    \bul{Won \textbf{three HPE hackathons}; two were selected for \textbf{HPE Fellow Review} and are being implemented across the business unit (an AI agent, its harness and a self-learning component). \textbf{Patent filing under way.}}
  \resumeItemListEnd

  \resumeSubheading{Hewlett Packard Enterprise, Storage}{San Jose, US}{Cloud Software Engineer}{Jun 2023 -- Feb 2026}
  \resumeItemListStart
    \bul{Built the \textbf{observability platform} for \textbf{30K devices} and \textbf{10K customers}: log-level root-cause analysis across hardware, OS, backend, UI and cloud cut resolution time from \textbf{30 to 3 minutes} and raised first-time fix success from \textbf{55\% to 85\%}.}
    \bul{Built monitoring dashboards and \textbf{Spark} log pipelines used by the teams supporting \textbf{10,000+ customers}, and visualized shipment and setup logs to find setup-time trends and anomalies.}
    \bul{Shipped \textbf{6 full-stack Peer Motion features} in Go and JavaScript across 3PAR, Primera, Alletra and Nimble arrays: Remote Copy support, the Nimble-to-Alletra iSCSI migration path, and domain, multi-link and VLAN tagging for migration flows. \textbf{15 tagged builds}, \textbf{47 defects} fixed across \textbf{5 release lines}.}
    \bul{Enabled cloud-based storage management for customer self-installation of new arrays, working with the QA and OS teams.}
    \bul{Built a \textbf{documentation agent} over a vector knowledge base that updates docs and API references on every pull request, saving about \textbf{5 engineer-hours a week per team}.}
    \bul{Proposed a \textbf{cross-business-unit observability and traceability framework} joining telemetry from Storage (Alletra), Networking (Aruba, Juniper) and Server platforms.}
    \bul{\textbf{2024 HPE hackathon winner} with an end-to-end RAG chatbot over PDFs; named \textbf{HPE Early Career Technologist}.}
  \resumeItemListEnd

  \resumeSubheading{The Luminosity Lab, Arizona State University}{Tempe, US}{Research Aide}{Feb 2022 -- May 2023}
  \resumeItemListStart
    \bul{Led the iOS team revamping the \href{https://thetouchpointsolution.com/}{TouchPoint} iOS and Android apps and helped the client make mobile product decisions.}
    \bul{Used HealthKit heart-rate-variability and sensor data to classify users into six stress personalities.}
    \bul{Designed a telehealth flow on the Zoom Meeting SDK, with Phoenix Children's Hospital, that lets doctors start remote visits and drive the patient's Android TV, saving \textbf{15 minutes} per visit.}
  \resumeItemListEnd

  \resumeSubheading{Aruba, Hewlett Packard Enterprise}{Tempe, US}{Software Engineering Intern}{May 2022 -- Aug 2022}
  \resumeItemListStart
    \bul{Built a test-automation tool that configures network topologies, saving an estimated \textbf{150k engineering hours a year}.}
    \bul{Presented the product to panels of senior engineers and cross-functional representatives.}
  \resumeItemListEnd

  \resumeSubheading{Webex by Cisco}{Bengaluru, India}{Technical Undergraduate Intern}{Jan 2021 -- Jul 2021}
  \resumeItemListStart
    \bul{Implemented call and meeting scheduling APIs for the \href{https://developer.webex.com/docs/sdks/ios}{Webex iOS SDK v3}.}
    \bul{Built \href{https://github.com/webex/webex-ios-sdk-example}{Kitchen Sink}, the sample iOS app that demonstrates the SDK's APIs.}
  \resumeItemListEnd
\resumeSubHeadingListEnd

%---------- SKILLS ----------
\section{Skills}
\resumeSubHeadingListStart
  \resumeSubItem{\textbullet\ Generative AI (GenAI)}{LLM agents, agentic workflows, multi-agent orchestration, MCP and tool calling, evals and LLM-as-a-judge, Retrieval-Augmented Generation (RAG), embeddings and vector search (pgvector), guardrails, prompt and context engineering, memory, tracing, LangGraph, vLLM, Qwen, llama.cpp, Ollama, Whisper}
  \resumeSubItem{\textbullet\ Languages}{Python, Go, TypeScript, JavaScript, Swift, Kotlin, C++, SQL, Shell}
  \resumeSubItem{\textbullet\ Infrastructure}{Distributed systems, Kubernetes, Helm, Docker, AWS (Lambda, SQS, EC2, S3, DynamoDB), Kafka, Spark, Trino, Postgres, Supabase, Firebase, Prometheus, Grafana, Humio, gRPC, REST, FastAPI, Node.js, Jenkins CI/CD, Git, Linux}
  \resumeSubItem{\textbullet\ Apps and UI}{React, Next.js, Tailwind, SwiftUI, UIKit, HealthKit, Android, Playwright}
  \resumeSubItem{\textbullet\ Product}{Problem scoping with users, success metrics and pass bars, zero-to-one delivery, cross-functional presentations}
\resumeSubHeadingListEnd

%---------- PROJECTS ----------
\section{Projects}
\resumeSubHeadingListStart
  \resumeSubItem{\textbullet\ Zuo, personal AI harness (2026)}{Job queue, budgets, project board and agent teams controlled from my phone; free NVIDIA endpoints, a PC GPU on vLLM and a laptop on llama.cpp as workers, Claude as planner. Voice capture with silence gating and Whisper; an iOS client; a UI agent team that turns design posts into web sections. Python, FastAPI, LangGraph, Docker.}
  \resumeSubItem{\textbullet\ Catan Arena, self-learning agents (2026)}{One learning agent plays full Catan against three frozen copies of itself; it graduates only when the 95\% Wilson interval of its win rate clears the 25\% null. Actor, watcher and consolidator loop with persistent memory.}
  \resumeSubItem{\textbullet\ AI Tutor (2026)}{Adaptive learning app for AI engineering: a level path, quizzes and a swipeable feed built from Reddit, Hacker News, arXiv and YouTube; Claude with an automatic NVIDIA Nemotron fallback.}
  \resumeSubItem{\textbullet\ Zzz, iOS and Android (2026)}{Gamified goal tracker with quests, points and rewards, live on the App Store. SwiftUI, Kotlin, Supabase.}
  \resumeSubItem{\textbullet\ Drone inspection simulation (2026)}{PX4 and Gazebo inspection drone with an agent capped to Raspberry Pi 5 resources: YOLOv8 plus a small vision-language model; llama.cpp 18x faster than Ollama on the same limits.}
  \resumeSubItem{\textbullet\ musaic (2026)}{Automated lo-fi music channel: Stable Audio, GPT-4 metadata, DALL-E and Imagen thumbnails, Veo backgrounds, FFmpeg encode and YouTube upload. \href{https://github.com/Fluff18/musaic}{GitHub}}
  \resumeSubItem{\textbullet\ Quantum slot machine (2025)}{Randomness from real qubits on IBM Quantum hardware; Qiskit and FastAPI backend, React frontend. \href{https://github.com/Fluff18/QuantumSlot}{GitHub}}
  \resumeSubItem{\textbullet\ Earlier (2020 to 2022)}{RideShare, a fault-tolerant database service on AWS with RabbitMQ and Zookeeper; face recognition on AWS Lambda and SQS (93\% accuracy) and on autoscaled EC2 (100 concurrent requests); handwriting recognition split across phones (30\% faster); Fetch Rewards recipe app in Swift.}
\resumeSubHeadingListEnd

%---------- RECOGNITION ----------
\section{Recognition and Leadership}
\resumeSubHeadingListStart
  \resumeSubItem{\textbullet\ HPE}{Three hackathon wins (2024, 2026); two selected for Fellow Review and BU-wide implementation; Early Career Technologist (2024); patent filing under way.}
  \resumeSubItem{\textbullet\ PES Center for Innovation and Entrepreneurship}{Led a team of six from idea to prototype of a smart attendance wristband; Best Project at Building a Lean Startup across 34 teams.}
  \resumeSubItem{\textbullet\ Stanford Continuing Studies}{Building Interpersonal Skills (COM 19), the T-group workshop modelled on Stanford GSB's Interpersonal Dynamics.}
\resumeSubHeadingListEnd

%---------- PUBLICATIONS ----------
\section{Publications}
\resumeSubHeadingListStart
  \bul{N. Raju, \textbf{R. M. Koushik}, H. Om, S. Kalambur. \textit{ViStA: Video Streaming and Analytics Benchmark.} IEEE International Symposium on Performance Analysis of Systems and Software (ISPASS), 2021.}
  \bul{\textbf{R. M. Koushik}, A. Perichiappan, H. Om, A. Banerji, S. Eswaran, P. Honnavalli. \textit{Generation of True Random Numbers using Entropy Sources Present within Portable Computers.} IEEE CONECCT, 2021.}
  \bul{R. Chandan, N. Pentapati, \textbf{R. M. Koushik}, R. Nagpal. \textit{A High Performance Pipelined Parallel Generative Adversarial Network (PipeGAN).} Lecture Notes in Networks and Systems, Springer, vol. 171, pp. 769--777, ICICSE 2020.}
\resumeSubHeadingListEnd

%---------- EDUCATION ----------
\section{Education}
\resumeSubHeadingListStart
  \resumeSubheading{Arizona State University}{Tempe, US}{Master of Science, Computer Science. GPA 3.59}{Aug 2021 -- May 2023}
  \resumeSubheading{PES University}{Bengaluru, India}{Bachelor of Technology, Computer Science and Engineering. GPA 8.2}{Aug 2017 -- Jul 2021}
\resumeSubHeadingListEnd

\end{document}
